\section{Method}
\label{sec:method}

\subsection{Machines}

Table~\ref{tab:machines} lists the five CPUs.
The three cloud types are AWS EC2 on-demand instances in us-east-1.
Each held-out set ran whole on one instance; four instances per type shared the sets (shards), and a fifth instance per type (the replica) repeated agent-ho and gcsn-ho.
The Linux machines report no frequency governor (the hypervisor controls frequency).
No arm was pinned to a core; the one-thread runs set a thread budget of one.
All Linux machines used Python~3.12.3, NumPy~2.5.4 on OpenBLAS~0.3.34 (BLIS~2.0 for one AMD variant), PyTorch~2.14.1 (CPU build), Node.js~24.18.0 and Rust~1.97.1; the Mac used Python~3.12.13 and Apple Accelerate.
Lock files pin every Python and Rust dependency.

\begin{table}[!htbp]
\centering\small
\caption{CPUs. Physical cores $\times$ hardware threads per core. ``jobs'' counts the held-out result jobs that passed the completeness check (shards; replica in parentheses).}
\label{tab:machines}
\resizebox{\linewidth}{!}{%
\begin{tabular}{llllr}
\toprule
name & CPU & cores $\times$ threads & vector units & jobs \\
\midrule
Intel & Xeon 6 6975P-C (Granite Rapids), c8i.4xlarge & $8\times2$ & AVX-512, AMX & \nm{count.intel} (\nm{count.intel-replica}) \\
AMD & EPYC 9R45 (Zen 5), c8a.4xlarge & $16\times1$ & AVX-512 & \nm{count.amd} (\nm{count.amd-replica}) \\
Graviton4 & Neoverse V2, c8g.4xlarge & $16\times1$ & NEON, SVE2 & \nm{count.arm} (\nm{count.arm-replica}) \\
M3 Max & Apple M3 Max, macOS 26.3.1 & 12 P + 4 E & NEON (+ AMX via Accelerate) & \nm{count.m3max} \\
EPYC & AMD EPYC (Zen 3), Hetzner CCX33 & \TODO{epyc} & AVX2 & \TODO{epyc} \\
\bottomrule
\end{tabular}}
\end{table}

\subsection{Arms and the frozen shortlist}

Each library ran in several configurations: API variant, input type (float32, bfloat16, float16; for NumPy and maxsim-cpu also ``int4-self'', where the rival decodes the int4 storage inside its timed call), and thread budget (library default, the physical or performance cores, and one).
On the development sets, a script chose per machine, library, input type and workload family the fastest two configurations and the fastest one-thread configuration.
This \emph{shortlist} was frozen with everything else.
Every table reports the frozen configuration; the fastest configuration on the held-out calls themselves appears only in the released per-machine summaries, labelled ``post hoc''.

We gave the rivals their best case.
Their inputs are decoded, normalized, cast and packed outside the timer (and counted in the memory measurement).
NumKong has no batch interface, so besides its Python interface we looped its C kernel over the documents in Rust (``native'' arms); NextPlaid's scorer was looped the same way, single-threaded and with a thread pool.
maxsim-cpu's x86 build returns wrong scores for queries over 32 tokens once a document has 64 tokens (Section~\ref{sec:correctness}); on x86 we therefore call it in 32-token query chunks, whose sums are exact.
Intel MKL NumPy did not install from Intel's package index and is missing.
maxsim-cpu has no Linux AArch64 wheel and is missing on Graviton4.

\subsection{Timing}

Each arm runs in its own process, because rival thread pools keep spinning after a call (OpenBLAS for up to about 100\,ms) and would slow the next arm.
Before each arm, the runner waits until machine CPU use, sampled over one second, is below 10\% (15\% on the Mac).
Per call, one untimed warm-up precedes the timed repetitions.
\begin{description}
\item[Warm] calls back to back; 5 timed repetitions per call.
\item[Pool-parked (5\,ms busy gap)] the main thread spins for 5\,ms before every call, so worker pools that park after a short idle time have parked; 3 timed repetitions. This mode does \emph{not} test cold caches or process start.
\item[One thread] the arms' one-thread configurations, warm.
\item[Parallel callers] see Section~\ref{sec:conc}.
\end{description}
Every arm runs $R = 5$ process repetitions $\times$ 2 rounds, the second round in reversed arm order, so that slow drift affects all arms alike.
A call's latency is the median over its 10 runs of the per-run median.
We keep every individual timing: ``per-request p95/p99'' pools all of them and is the only tail statistic we report.

\paragraph{Ratios and uncertainty.}
A speed ratio is the geometric mean over calls of (rival time / \sweet time); above 1 means \sweet is faster.
We report two uncertainty numbers apart.
The \emph{workload interval} is a 95\% cluster-bootstrap interval (10,000 resamples, seed 42) that resamples repositories for agent calls (11 clusters), queries for GCSN and BEIR, and calls for sweep cells.
The \emph{execution spread} is the range of the ratio over the 5 process repetitions; the replica instance adds an instance-to-instance spread (Section~\ref{sec:spread}).

\paragraph{Noise check and $R$.}
On development data, an A/A run (the same build as two arms, 5 process repetitions) gave per machine the coefficient of variation $cv$ of the set-level geometric mean and the noise band $n = \max(1\%, 2\,cv)$: Intel \nm{noise.intel}\%, AMD \nm{noise.amd}\%, Graviton4 \nm{noise.arm}\%, M3 Max \nm{noise.m3max}\%, EPYC \nm{noise.epyc}\%.
The rule $R = \operatorname{clamp}(\lceil (cv/0.01)^2\rceil, 3, 8)$ gave 3 or 4. We used $R = \max(5, \text{noise rule}) = 5$ on all machines (an early plan with $R = 3$ was corrected before any held-out run).

\subsection{Correctness on every call}

Every timed run also checks every score.
Per arm we report the largest absolute error against reference (a) of its own contract, the input-rounding effect, and the agreement of its top-ranked document with the contract-S reference (``top-1 = S'').
For the BEIR sets every arm stores all its scores, so each implementation gets its own quality curve.

\subsection{Development, freeze and held-out runs}
\label{sec:freeze}

All tuning used development data only.
The last kernel change (the AVX-512 block) was kept on development data under a pre-set rule: at least 3\% faster and above the noise band in two rounds on the target machine, no slowdown above the band elsewhere, and bit-identical scores.
A second candidate (a thread-count rule by call size) failed that rule and was dropped.

Then we froze, in one file, the code state, every adapter, the shortlists, the analysis scripts, the dataset definitions, the contracts, $R$, the quality margin $\delta$, the four primary quality tests, the parallel-caller rates, and a rule for a stronger claim.
Three commits define the freeze: \texttt{a4317ba9} (the frozen code), \texttt{ea84951f} (the freeze file re-made after a replica naming fix; the held-out bundle came from it) and \texttt{daf54122} (a memory fix in the parallel-caller harness with its new freeze entry, applied to every machine before any reached a parallel-caller group).
A script checks the SHA-256 of every frozen file before use.
Analysis code written after the freeze lives in a separate folder and is named so here.

\paragraph{Held-out discipline.}
We read held-out results only as aggregates.
No per-query or per-call held-out row was printed or inspected, and nothing was tuned after a held-out result was seen.
